\section{Introduction}

\begin{figure}[!t]
  \centering
  \includesvg[width=1.0\linewidth]{img/dataset_creation_2} %
  \caption{
    \DatasetName dataset creation (left) and experimental setup (right). 
  }
  \label{fig:dataset_creation}
\end{figure}

Securing intellectual property is a long and costly process that requires both deep technical knowledge and expertise in patent law.
This motivates the use of technology to boost patent attorney productivity.
Natural language processing (NLP) already assists prior art search \cite{shalabyPatentRetrievalLiterature2019,stamatisEndEndNeural2022,pujariMultitaskApproachNeural2021} and patent landscaping \citep{CHOI2022121413,pujariThreeRealWorldDatasets2022}.
Research on Large Language Models (LLMs) in the patent domain has recently gained momentum \cite{shomeeComprehensiveSurveyAIbased2024,jiangArtificialIntelligenceExploring2024,casolaSummarizationSimplificationGeneration2022,wangIPEvalBilingualIntellectual2024}, but patent drafting remains a largely manual task.

A patent typically consists of claims, which define the invention and
the legally relevant scope of protection, and a description, which provides technical details in sections like \textit{Field of the Invention}, \textit{Background}, \textit{Summary}, and \textit{Detailed Description}.
Prior work has primarily focused on generating abstracts and claims \cite{Hamborg2017Autom-41994,leePatentTransformer2ControllingPatent2020,christofidellisPGTPromptBased2022,leeEvaluatingGenerativePatent2023,zuoPatentEvalUnderstandingErrors2024,leeInstructPatentGPTTrainingPatent2024,baiPatentGPTLargeLanguage2024}.
The description makes up over 90\% of the document,
\footnote{Measured on our dataset using the Llama-3 tokenizer: 0.7\% abstract, 91.8\% description, 7.5\% claims.}
implying that large productivity gains are expected from writing support for this section.
Yet, generating them remains a significant challenge for LLMs due to their length, technical complexity, and specialized language
\cite{wangPatentformerNovelMethod2024,wangAutoPatentMultiAgentFramework2024}.
Existing work on automatic patent generation suffers from a number of shortcomings, such as ill-posed task setups, lack of open benchmarks, and disregard for patent descriptions \citep{jiangArtificialIntelligenceExploring2024}.
In this paper, we tackle all of these issues: we propose a novel setup, evaluation metrics, and outline-guided models for description generation %
in a realistic setting using open and human-created data.

Inventors typically submit invention reports (IRs), which patent attorneys formalize into patent applications.
In many research labs, it is common to use a pre-publication paper as invention report, which leads to so-called \textit{patent-paper pairs} (PPPs, \citet{murrayInnovationCoevolutionScientific2002}), an unrecognized treasure for AI research. %
To facilitate the study of LLMs on patent drafting, we create \DatasetName, a new benchmark of 1.8k carefully identified PPPs and patent outlines.
We develop and validate a method for reliably matching patents and papers describing the same invention (see \Cref{fig:dataset_creation}).
For LLM-supported patent drafting, we envision a practical setting in which the attorney, given a paper, %
provides an outline for the patent.
This outline acts as a flexible mechanism to control the document structure and content while keeping manual effort low.

A major challenge of the proposed task is document length: patent descriptions in \DatasetName are on average 18k tokens long, some exceeding 180k tokens.
While current LLMs increasingly support long context windows, they struggle to generate similarly long outputs \cite{liuLongGenBenchLongcontextGeneration2024,baiLongWriterUnleashing102024,wuLongGenBenchBenchmarkingLongForm2025,yeLongProcBenchmarkingLongContext2025}.
As a remedy, we propose \textit{chunk-based outline-guided patent generation} (\MethodName), which effectively generates long patent documents in chunks.

Evaluating generated patents poses significant challenges due to their length, their technical complexity, and the high cost of manual evaluation.
No standard evaluation metrics are established in the literature, and prior work \cite{wangAutoPatentMultiAgentFramework2024,wangPatentformerNovelMethod2024} commonly resorts to standard text similarity metrics which do not work well on long documents \cite{lattimerFastAccurateFactual2023,queHelloBenchEvaluatingLong2024}. 
We adapt a suite of metrics based on natural language inference (NLI) and authorship attribution to the specific case of evaluating factual correctness, coverage, and language style for long-form patent generation on \DatasetName.
Our main contributions are:
\begin{enumerate}[label={(\arabic*)},itemsep=0mm,wide,labelindent=0pt,topsep=0pt,partopsep=0ex,parsep=0ex]
  \item We create and release \DatasetName, a new benchmark for patent drafting based on open data that closely aligns with real-world settings.

  \item We derive a comprehensive suite of automatic evaluation metrics for evaluating generated patents.

  \item We propose a chunk-based outline-guided patent generation approach with effectively controllable output length. 
  Our method increases coverage considerably while keeping factuality high.

  \item We conduct extensive evaluations finding that state-of-the-art LLMs can effectively use information from the papers, but still struggle to provide the necessary level of detail,
  and that fine-tuning leads to much higher stylistic similarity with patents, but substantially decreased factuality.

  \item Our human evaluation confirms these findings and indicates promising potential for productivity gains of patent attorneys.
\end{enumerate}

